
\documentclass{article}
\usepackage{colortbl}
\usepackage{makecell}
\usepackage{multirow}
\usepackage{supertabular}

\begin{document}

\newcounter{utterance}

\centering \large Interaction Transcript for game `codenames', experiment `opponent\_none', episode 1 with glm{-}5.2.
\vspace{24pt}

{ \footnotesize  \setcounter{utterance}{1}
\setlength{\tabcolsep}{0pt}
\begin{supertabular}{c@{$\;$}|p{.15\linewidth}@{}p{.15\linewidth}p{.15\linewidth}p{.15\linewidth}p{.15\linewidth}p{.15\linewidth}}
    \# & $\;$A & \multicolumn{4}{c}{Game Master} & $\;\:$B\\
    \hline

    \theutterance \stepcounter{utterance}  
    & & & \multicolumn{2}{p{0.3\linewidth}}{
        \cellcolor[rgb]{0.9,0.9,0.9}{
            \makecell[{{p{\linewidth}}}]{
                \texttt{\tiny{[GM$|$GM]}}
                {'hidden': {'team': ['plan', 'antarktis', 'bomb', 'citron', 'peking', 'oliv', 'pilot', 'pund', 'mikroskop'], 'innocent': ['triangel', 'mössa', 'modell', 'ändra', 'kort', 'lager', 'stat'], 'opponent': ['hand', 'kraft', 'fågel fenix', 'tokyo', 'våg', 'kung', 'elev', 'obligation'], 'assassin': ['papper']}, 'revealed': {'team': {'team': [], 'innocent': [], 'opponent': [], 'assassin': []}, 'opponent': {'team': [], 'innocent': [], 'opponent': [], 'assassin': []}}}
            }
        }
    }
    & & \\ \\

    \theutterance \stepcounter{utterance}  
    & & \multicolumn{4}{p{0.6\linewidth}}{
        \cellcolor[rgb]{0.9,0.9,0.9}{
            \makecell[{{p{\linewidth}}}]{
                \texttt{\tiny{[P1$\langle$GM]}}
                \texttt{Ge en ledtråd bestående av ett enda ord som relaterar till några av ditt lags ord från listan nedan. Du kan välja hur många och vilka ord din ledtråd ska relatera till.} \\
\texttt{Målet är att få din partner att gissa alla ditt lags ord baserat på de ledtrådar du ger dem, innan det andra laget har avslöjat alla sina. Om din partner gissar ett "motståndare"{-}ord eller ett "distraktions"{-}ord avslutas deras gissningstur och motståndarlaget får göra en tur. Om din partner gissar ett "död"{-}ord slutar spelet och ditt lag förlorar omedelbart.} \\
\texttt{Ju fler ord din ledtråd relaterar till, desto snabbare kan du vinna spelet, men desto svårare kan associationerna vara att gissa för din partner, så välj ditt ledord klokt. Ledordet måste vara semantiskt relaterat till målorden, det får inte vara ett av orden i listorna eller innehålla delar av dem.} \\
\\ 
\texttt{Ge alltid din ledtråd bestående av ett enda ord och din kommaseparerade lista över relaterade målord i följande format och gör dina svar så korta som möjligt, inkludera aldrig någon annan text än vad som krävs i detta format:} \\
\\ 
\texttt{LEDTRÅD: <ORD>} \\
\texttt{MÅL: <MÅL>} \\
\\ 
\texttt{Dina lagord är: plan, antarktis, bomb, citron, peking, oliv, pilot, pund, mikroskop.} \\
\texttt{Din motståndares lagord är: hand, kraft, fågel fenix, tokyo, våg, kung, elev, obligation.} \\
\texttt{Distraktionsord är: triangel, mössa, modell, ändra, kort, lager, stat.} \\
\texttt{Dödorden är: papper.} \\
            }
        }
    }
    & \\ \\

    \theutterance \stepcounter{utterance}  
    & \multicolumn{4}{p{0.6\linewidth}}{
        \cellcolor[rgb]{0.9,0.9,0.9}{
            \makecell[{{p{\linewidth}}}]{
                \texttt{\tiny{[P1$\rangle$GM]}}
                \texttt{LEDTRÅD: citrus} \\
\texttt{MÅL: citron, oliv} \\
            }
        }
    }
    & & \\ \\

    \theutterance \stepcounter{utterance}  
    & & & \multicolumn{2}{p{0.3\linewidth}}{
        \cellcolor[rgb]{0.9,0.9,0.9}{
            \makecell[{{p{\linewidth}}}]{
                \texttt{\tiny{[GM$|$GM]}}
                \texttt{citrus} \\
            }
        }
    }
    & & \\ \\

    \theutterance \stepcounter{utterance}  
    & & & \multicolumn{2}{p{0.3\linewidth}}{
        \cellcolor[rgb]{0.9,0.9,0.9}{
            \makecell[{{p{\linewidth}}}]{
                \texttt{\tiny{[GM$|$GM]}}
                ['citron', 'oliv']
            }
        }
    }
    & & \\ \\

    \theutterance \stepcounter{utterance}  
    & & & \multicolumn{2}{p{0.3\linewidth}}{
        \cellcolor[rgb]{0.9,0.9,0.9}{
            \makecell[{{p{\linewidth}}}]{
                \texttt{\tiny{[GM$|$GM]}}
                {'word': 'citron', 'assignment': 'team'}
            }
        }
    }
    & & \\ \\

    \theutterance \stepcounter{utterance}  
    & & & \multicolumn{2}{p{0.3\linewidth}}{
        \cellcolor[rgb]{0.9,0.9,0.9}{
            \makecell[{{p{\linewidth}}}]{
                \texttt{\tiny{[GM$|$GM]}}
                {'word': 'oliv', 'assignment': 'team'}
            }
        }
    }
    & & \\ \\

    \theutterance \stepcounter{utterance}  
    & & \multicolumn{4}{p{0.6\linewidth}}{
        \cellcolor[rgb]{0.9,0.9,0.9}{
            \makecell[{{p{\linewidth}}}]{
                \texttt{\tiny{[GM$\rangle$P2]}}
                \texttt{Ge en kommaseparerad lista på upp till 2 ord från följande lista som bäst relaterar till eller är närmast associerade med ordet 'citrus'. Börja alltid din lista med gissningar med GISSNING: och ta inte med någon annan text i ditt svar.} \\
\\ 
\texttt{kraft, antarktis, hand, lager, pund, plan, oliv, bomb, ändra, fågel fenix, obligation, modell, kung, mössa, citron, våg, tokyo, peking, triangel, papper, pilot, kort, stat, mikroskop, elev} \\
            }
        }
    }
    & \\ \\

    \theutterance \stepcounter{utterance}  
    & & & \multicolumn{4}{p{0.6\linewidth}}{
        \cellcolor[rgb]{0.9,0.9,0.9}{
            \makecell[{{p{\linewidth}}}]{
                \texttt{\tiny{[GM$\langle$P2]}}
                \texttt{GISSNING: citron, oliv} \\
            }
        }
    }
     \\ \\

    \theutterance \stepcounter{utterance}  
    & & & \multicolumn{2}{p{0.3\linewidth}}{
        \cellcolor[rgb]{0.9,0.9,0.9}{
            \makecell[{{p{\linewidth}}}]{
                \texttt{\tiny{[GM$|$GM]}}
                ['citron', 'oliv']
            }
        }
    }
    & & \\ \\

    \theutterance \stepcounter{utterance}  
    & & & \multicolumn{2}{p{0.3\linewidth}}{
        \cellcolor[rgb]{0.9,0.9,0.9}{
            \makecell[{{p{\linewidth}}}]{
                \texttt{\tiny{[GM$|$GM]}}
                {'word': 'citron', 'assignment': 'team'}
            }
        }
    }
    & & \\ \\

    \theutterance \stepcounter{utterance}  
    & & & \multicolumn{2}{p{0.3\linewidth}}{
        \cellcolor[rgb]{0.9,0.9,0.9}{
            \makecell[{{p{\linewidth}}}]{
                \texttt{\tiny{[GM$|$GM]}}
                {'word': 'citron', 'assignment': 'team'}
            }
        }
    }
    & & \\ \\

    \theutterance \stepcounter{utterance}  
    & & & \multicolumn{2}{p{0.3\linewidth}}{
        \cellcolor[rgb]{0.9,0.9,0.9}{
            \makecell[{{p{\linewidth}}}]{
                \texttt{\tiny{[GM$|$GM]}}
                {'word': 'oliv', 'assignment': 'team'}
            }
        }
    }
    & & \\ \\

    \theutterance \stepcounter{utterance}  
    & & & \multicolumn{2}{p{0.3\linewidth}}{
        \cellcolor[rgb]{0.9,0.9,0.9}{
            \makecell[{{p{\linewidth}}}]{
                \texttt{\tiny{[GM$|$GM]}}
                {'word': 'oliv', 'assignment': 'team'}
            }
        }
    }
    & & \\ \\

    \theutterance \stepcounter{utterance}  
    & & & \multicolumn{2}{p{0.3\linewidth}}{
        \cellcolor[rgb]{0.9,0.9,0.9}{
            \makecell[{{p{\linewidth}}}]{
                \texttt{\tiny{[GM$|$GM]}}
                {'hidden': {'team': ['plan', 'antarktis', 'bomb', 'peking', 'pilot', 'pund', 'mikroskop'], 'innocent': ['triangel', 'mössa', 'modell', 'ändra', 'kort', 'lager', 'stat'], 'opponent': ['hand', 'kraft', 'fågel fenix', 'tokyo', 'våg', 'kung', 'elev', 'obligation'], 'assassin': ['papper']}, 'revealed': {'team': {'team': ['citron', 'oliv'], 'innocent': [], 'opponent': [], 'assassin': []}, 'opponent': {'team': [], 'innocent': [], 'opponent': [], 'assassin': []}}}
            }
        }
    }
    & & \\ \\

    \theutterance \stepcounter{utterance}  
    & & \multicolumn{4}{p{0.6\linewidth}}{
        \cellcolor[rgb]{0.9,0.9,0.9}{
            \makecell[{{p{\linewidth}}}]{
                \texttt{\tiny{[P1$\langle$GM]}}
                \texttt{Orden citron, oliv gissades korrekt. Din lagkamrats tur slutade där.} \\
\texttt{Ge nu en ny ledtråd som relaterar till några av dina återstående lagord och en lista med relaterade målord. Kom ihåg att börja din ledtråd med 'LEDTRÅD: ', skriv en ny rad och börja din kommaseparerade lista med målord med 'MÅL: '. Observera: vissa ord har tagits bort från listorna jämfört med tidigare förfrågningar.} \\
\\ 
\texttt{Dina återstående lagord är: plan, antarktis, bomb, peking, pilot, pund, mikroskop.} \\
\texttt{Återstående ord för din motståndare är: hand, kraft, fågel fenix, tokyo, våg, kung, elev, obligation.} \\
\texttt{Återstående distraherande ord är: triangel, mössa, modell, ändra, kort, lager, stat.} \\
\texttt{Återstående dödande ord är: papper.} \\
            }
        }
    }
    & \\ \\

    \theutterance \stepcounter{utterance}  
    & \multicolumn{4}{p{0.6\linewidth}}{
        \cellcolor[rgb]{0.9,0.9,0.9}{
            \makecell[{{p{\linewidth}}}]{
                \texttt{\tiny{[P1$\rangle$GM]}}
                \texttt{LEDTRÅD: kina} \\
\texttt{MÅL: peking} \\
            }
        }
    }
    & & \\ \\

    \theutterance \stepcounter{utterance}  
    & & & \multicolumn{2}{p{0.3\linewidth}}{
        \cellcolor[rgb]{0.9,0.9,0.9}{
            \makecell[{{p{\linewidth}}}]{
                \texttt{\tiny{[GM$|$GM]}}
                \texttt{kina} \\
            }
        }
    }
    & & \\ \\

    \theutterance \stepcounter{utterance}  
    & & & \multicolumn{2}{p{0.3\linewidth}}{
        \cellcolor[rgb]{0.9,0.9,0.9}{
            \makecell[{{p{\linewidth}}}]{
                \texttt{\tiny{[GM$|$GM]}}
                ['peking']
            }
        }
    }
    & & \\ \\

    \theutterance \stepcounter{utterance}  
    & & & \multicolumn{2}{p{0.3\linewidth}}{
        \cellcolor[rgb]{0.9,0.9,0.9}{
            \makecell[{{p{\linewidth}}}]{
                \texttt{\tiny{[GM$|$GM]}}
                {'word': 'peking', 'assignment': 'team'}
            }
        }
    }
    & & \\ \\

    \theutterance \stepcounter{utterance}  
    & & \multicolumn{4}{p{0.6\linewidth}}{
        \cellcolor[rgb]{0.9,0.9,0.9}{
            \makecell[{{p{\linewidth}}}]{
                \texttt{\tiny{[GM$\rangle$P2]}}
                \texttt{Orden citron, oliv gissades korrekt. Din tur slutade där.} \\
\texttt{Ange nu ytterligare en kommaseparerad lista med minst 1 och upp till 1 ord från följande lista med ord som bäst relaterar till eller är närmast associerade med ordet 'kina'. Kom ihåg att börja ditt svar med GISSNING:. Observera: vissa ord har tagits bort från listan jämfört med tidigare förfrågningar.} \\
\\ 
\texttt{kraft, antarktis, hand, lager, pund, plan, bomb, ändra, fågel fenix, obligation, modell, kung, mössa, våg, tokyo, peking, triangel, papper, pilot, kort, stat, mikroskop, elev} \\
            }
        }
    }
    & \\ \\

    \theutterance \stepcounter{utterance}  
    & & & \multicolumn{4}{p{0.6\linewidth}}{
        \cellcolor[rgb]{0.9,0.9,0.9}{
            \makecell[{{p{\linewidth}}}]{
                \texttt{\tiny{[GM$\langle$P2]}}
                \texttt{GISSNING: peking} \\
            }
        }
    }
     \\ \\

    \theutterance \stepcounter{utterance}  
    & & & \multicolumn{2}{p{0.3\linewidth}}{
        \cellcolor[rgb]{0.9,0.9,0.9}{
            \makecell[{{p{\linewidth}}}]{
                \texttt{\tiny{[GM$|$GM]}}
                ['peking']
            }
        }
    }
    & & \\ \\

    \theutterance \stepcounter{utterance}  
    & & & \multicolumn{2}{p{0.3\linewidth}}{
        \cellcolor[rgb]{0.9,0.9,0.9}{
            \makecell[{{p{\linewidth}}}]{
                \texttt{\tiny{[GM$|$GM]}}
                {'word': 'peking', 'assignment': 'team'}
            }
        }
    }
    & & \\ \\

    \theutterance \stepcounter{utterance}  
    & & & \multicolumn{2}{p{0.3\linewidth}}{
        \cellcolor[rgb]{0.9,0.9,0.9}{
            \makecell[{{p{\linewidth}}}]{
                \texttt{\tiny{[GM$|$GM]}}
                {'word': 'peking', 'assignment': 'team'}
            }
        }
    }
    & & \\ \\

    \theutterance \stepcounter{utterance}  
    & & & \multicolumn{2}{p{0.3\linewidth}}{
        \cellcolor[rgb]{0.9,0.9,0.9}{
            \makecell[{{p{\linewidth}}}]{
                \texttt{\tiny{[GM$|$GM]}}
                {'hidden': {'team': ['plan', 'antarktis', 'bomb', 'pilot', 'pund', 'mikroskop'], 'innocent': ['triangel', 'mössa', 'modell', 'ändra', 'kort', 'lager', 'stat'], 'opponent': ['hand', 'kraft', 'fågel fenix', 'tokyo', 'våg', 'kung', 'elev', 'obligation'], 'assassin': ['papper']}, 'revealed': {'team': {'team': ['citron', 'oliv', 'peking'], 'innocent': [], 'opponent': [], 'assassin': []}, 'opponent': {'team': [], 'innocent': [], 'opponent': [], 'assassin': []}}}
            }
        }
    }
    & & \\ \\

    \theutterance \stepcounter{utterance}  
    & & \multicolumn{4}{p{0.6\linewidth}}{
        \cellcolor[rgb]{0.9,0.9,0.9}{
            \makecell[{{p{\linewidth}}}]{
                \texttt{\tiny{[P1$\langle$GM]}}
                \texttt{Ordet peking gissades korrekt. Din lagkamrats tur slutade där.} \\
\texttt{Ge nu en ny ledtråd som relaterar till några av dina återstående lagord och en lista med relaterade målord. Kom ihåg att börja din ledtråd med 'LEDTRÅD: ', skriv en ny rad och börja din kommaseparerade lista med målord med 'MÅL: '. Observera: vissa ord har tagits bort från listorna jämfört med tidigare förfrågningar.} \\
\\ 
\texttt{Dina återstående lagord är: plan, antarktis, bomb, pilot, pund, mikroskop.} \\
\texttt{Återstående ord för din motståndare är: hand, kraft, fågel fenix, tokyo, våg, kung, elev, obligation.} \\
\texttt{Återstående distraherande ord är: triangel, mössa, modell, ändra, kort, lager, stat.} \\
\texttt{Återstående dödande ord är: papper.} \\
            }
        }
    }
    & \\ \\

    \theutterance \stepcounter{utterance}  
    & \multicolumn{4}{p{0.6\linewidth}}{
        \cellcolor[rgb]{0.9,0.9,0.9}{
            \makecell[{{p{\linewidth}}}]{
                \texttt{\tiny{[P1$\rangle$GM]}}
                \texttt{LEDTRÅD: flygning} \\
\texttt{MÅL: pilot, plan, bomb} \\
            }
        }
    }
    & & \\ \\

    \theutterance \stepcounter{utterance}  
    & & & \multicolumn{2}{p{0.3\linewidth}}{
        \cellcolor[rgb]{0.9,0.9,0.9}{
            \makecell[{{p{\linewidth}}}]{
                \texttt{\tiny{[GM$|$GM]}}
                \texttt{flygning} \\
            }
        }
    }
    & & \\ \\

    \theutterance \stepcounter{utterance}  
    & & & \multicolumn{2}{p{0.3\linewidth}}{
        \cellcolor[rgb]{0.9,0.9,0.9}{
            \makecell[{{p{\linewidth}}}]{
                \texttt{\tiny{[GM$|$GM]}}
                ['pilot', 'plan', 'bomb']
            }
        }
    }
    & & \\ \\

    \theutterance \stepcounter{utterance}  
    & & & \multicolumn{2}{p{0.3\linewidth}}{
        \cellcolor[rgb]{0.9,0.9,0.9}{
            \makecell[{{p{\linewidth}}}]{
                \texttt{\tiny{[GM$|$GM]}}
                {'word': 'pilot', 'assignment': 'team'}
            }
        }
    }
    & & \\ \\

    \theutterance \stepcounter{utterance}  
    & & & \multicolumn{2}{p{0.3\linewidth}}{
        \cellcolor[rgb]{0.9,0.9,0.9}{
            \makecell[{{p{\linewidth}}}]{
                \texttt{\tiny{[GM$|$GM]}}
                {'word': 'plan', 'assignment': 'team'}
            }
        }
    }
    & & \\ \\

    \theutterance \stepcounter{utterance}  
    & & & \multicolumn{2}{p{0.3\linewidth}}{
        \cellcolor[rgb]{0.9,0.9,0.9}{
            \makecell[{{p{\linewidth}}}]{
                \texttt{\tiny{[GM$|$GM]}}
                {'word': 'bomb', 'assignment': 'team'}
            }
        }
    }
    & & \\ \\

    \theutterance \stepcounter{utterance}  
    & & \multicolumn{4}{p{0.6\linewidth}}{
        \cellcolor[rgb]{0.9,0.9,0.9}{
            \makecell[{{p{\linewidth}}}]{
                \texttt{\tiny{[GM$\rangle$P2]}}
                \texttt{Ordet peking gissades korrekt. Din tur slutade där.} \\
\texttt{Ange nu ytterligare en kommaseparerad lista med minst 1 och upp till 3 ord från följande lista med ord som bäst relaterar till eller är närmast associerade med ordet 'flygning'. Kom ihåg att börja ditt svar med GISSNING:. Observera: vissa ord har tagits bort från listan jämfört med tidigare förfrågningar.} \\
\\ 
\texttt{kraft, antarktis, hand, lager, pund, plan, bomb, ändra, fågel fenix, obligation, modell, kung, mössa, våg, tokyo, triangel, papper, pilot, kort, stat, mikroskop, elev} \\
            }
        }
    }
    & \\ \\

    \theutterance \stepcounter{utterance}  
    & & & \multicolumn{4}{p{0.6\linewidth}}{
        \cellcolor[rgb]{0.9,0.9,0.9}{
            \makecell[{{p{\linewidth}}}]{
                \texttt{\tiny{[GM$\langle$P2]}}
                \texttt{GISSNING: pilot, fågel fenix, våg} \\
            }
        }
    }
     \\ \\

    \theutterance \stepcounter{utterance}  
    & & & \multicolumn{2}{p{0.3\linewidth}}{
        \cellcolor[rgb]{0.9,0.9,0.9}{
            \makecell[{{p{\linewidth}}}]{
                \texttt{\tiny{[GM$|$GM]}}
                ['pilot', 'fågel fenix', 'våg']
            }
        }
    }
    & & \\ \\

    \theutterance \stepcounter{utterance}  
    & & & \multicolumn{2}{p{0.3\linewidth}}{
        \cellcolor[rgb]{0.9,0.9,0.9}{
            \makecell[{{p{\linewidth}}}]{
                \texttt{\tiny{[GM$|$GM]}}
                {'word': 'pilot', 'assignment': 'team'}
            }
        }
    }
    & & \\ \\

    \theutterance \stepcounter{utterance}  
    & & & \multicolumn{2}{p{0.3\linewidth}}{
        \cellcolor[rgb]{0.9,0.9,0.9}{
            \makecell[{{p{\linewidth}}}]{
                \texttt{\tiny{[GM$|$GM]}}
                {'word': 'pilot', 'assignment': 'team'}
            }
        }
    }
    & & \\ \\

    \theutterance \stepcounter{utterance}  
    & & & \multicolumn{2}{p{0.3\linewidth}}{
        \cellcolor[rgb]{0.9,0.9,0.9}{
            \makecell[{{p{\linewidth}}}]{
                \texttt{\tiny{[GM$|$GM]}}
                {'word': 'fågel fenix', 'assignment': 'opponent'}
            }
        }
    }
    & & \\ \\

    \theutterance \stepcounter{utterance}  
    & & & \multicolumn{2}{p{0.3\linewidth}}{
        \cellcolor[rgb]{0.9,0.9,0.9}{
            \makecell[{{p{\linewidth}}}]{
                \texttt{\tiny{[GM$|$GM]}}
                \texttt{fågel fenix} \\
            }
        }
    }
    & & \\ \\

    \theutterance \stepcounter{utterance}  
    & & & \multicolumn{2}{p{0.3\linewidth}}{
        \cellcolor[rgb]{0.9,0.9,0.9}{
            \makecell[{{p{\linewidth}}}]{
                \texttt{\tiny{[GM$|$GM]}}
                {'hidden': {'team': ['plan', 'antarktis', 'bomb', 'pund', 'mikroskop'], 'innocent': ['triangel', 'mössa', 'modell', 'ändra', 'kort', 'lager', 'stat'], 'opponent': ['hand', 'kraft', 'tokyo', 'våg', 'kung', 'elev', 'obligation'], 'assassin': ['papper']}, 'revealed': {'team': {'team': ['citron', 'oliv', 'peking', 'pilot'], 'innocent': [], 'opponent': ['fågel fenix'], 'assassin': []}, 'opponent': {'team': [], 'innocent': [], 'opponent': [], 'assassin': []}}}
            }
        }
    }
    & & \\ \\

    \theutterance \stepcounter{utterance}  
    & & \multicolumn{4}{p{0.6\linewidth}}{
        \cellcolor[rgb]{0.9,0.9,0.9}{
            \makecell[{{p{\linewidth}}}]{
                \texttt{\tiny{[P1$\langle$GM]}}
                \texttt{Ordet pilot gissades korrekt. Ordet fågel fenix gissades, men det är ett motståndare ord. Din lagkamrats tur slutade där.} \\
\texttt{Ge nu en ny ledtråd som relaterar till några av dina återstående lagord och en lista med relaterade målord. Kom ihåg att börja din ledtråd med 'LEDTRÅD: ', skriv en ny rad och börja din kommaseparerade lista med målord med 'MÅL: '. Observera: vissa ord har tagits bort från listorna jämfört med tidigare förfrågningar.} \\
\\ 
\texttt{Dina återstående lagord är: plan, antarktis, bomb, pund, mikroskop.} \\
\texttt{Återstående ord för din motståndare är: hand, kraft, tokyo, våg, kung, elev, obligation.} \\
\texttt{Återstående distraherande ord är: triangel, mössa, modell, ändra, kort, lager, stat.} \\
\texttt{Återstående dödande ord är: papper.} \\
            }
        }
    }
    & \\ \\

    \theutterance \stepcounter{utterance}  
    & \multicolumn{4}{p{0.6\linewidth}}{
        \cellcolor[rgb]{0.9,0.9,0.9}{
            \makecell[{{p{\linewidth}}}]{
                \texttt{\tiny{[P1$\rangle$GM]}}
                \texttt{LEDTRÅD: vetenskap} \\
\texttt{MÅL: mikroskop, plan} \\
            }
        }
    }
    & & \\ \\

    \theutterance \stepcounter{utterance}  
    & & & \multicolumn{2}{p{0.3\linewidth}}{
        \cellcolor[rgb]{0.9,0.9,0.9}{
            \makecell[{{p{\linewidth}}}]{
                \texttt{\tiny{[GM$|$GM]}}
                \texttt{vetenskap} \\
            }
        }
    }
    & & \\ \\

    \theutterance \stepcounter{utterance}  
    & & & \multicolumn{2}{p{0.3\linewidth}}{
        \cellcolor[rgb]{0.9,0.9,0.9}{
            \makecell[{{p{\linewidth}}}]{
                \texttt{\tiny{[GM$|$GM]}}
                ['mikroskop', 'plan']
            }
        }
    }
    & & \\ \\

    \theutterance \stepcounter{utterance}  
    & & & \multicolumn{2}{p{0.3\linewidth}}{
        \cellcolor[rgb]{0.9,0.9,0.9}{
            \makecell[{{p{\linewidth}}}]{
                \texttt{\tiny{[GM$|$GM]}}
                {'word': 'mikroskop', 'assignment': 'team'}
            }
        }
    }
    & & \\ \\

    \theutterance \stepcounter{utterance}  
    & & & \multicolumn{2}{p{0.3\linewidth}}{
        \cellcolor[rgb]{0.9,0.9,0.9}{
            \makecell[{{p{\linewidth}}}]{
                \texttt{\tiny{[GM$|$GM]}}
                {'word': 'plan', 'assignment': 'team'}
            }
        }
    }
    & & \\ \\

    \theutterance \stepcounter{utterance}  
    & & \multicolumn{4}{p{0.6\linewidth}}{
        \cellcolor[rgb]{0.9,0.9,0.9}{
            \makecell[{{p{\linewidth}}}]{
                \texttt{\tiny{[GM$\rangle$P2]}}
                \texttt{Ordet pilot gissades korrekt. Ordet fågel fenix gissades, men det är ett motståndare ord. Din tur slutade där.} \\
\texttt{Ange nu ytterligare en kommaseparerad lista med minst 1 och upp till 2 ord från följande lista med ord som bäst relaterar till eller är närmast associerade med ordet 'vetenskap'. Kom ihåg att börja ditt svar med GISSNING:. Observera: vissa ord har tagits bort från listan jämfört med tidigare förfrågningar.} \\
\\ 
\texttt{kraft, antarktis, hand, lager, pund, plan, bomb, ändra, obligation, modell, kung, mössa, våg, tokyo, triangel, papper, kort, stat, mikroskop, elev} \\
            }
        }
    }
    & \\ \\

    \theutterance \stepcounter{utterance}  
    & & & \multicolumn{4}{p{0.6\linewidth}}{
        \cellcolor[rgb]{0.9,0.9,0.9}{
            \makecell[{{p{\linewidth}}}]{
                \texttt{\tiny{[GM$\langle$P2]}}
                \texttt{GISSNING: mikroskop, laboratorium} \\
            }
        }
    }
     \\ \\

    \theutterance \stepcounter{utterance}  
    & & & \multicolumn{2}{p{0.3\linewidth}}{
        \cellcolor[rgb]{0.9,0.9,0.9}{
            \makecell[{{p{\linewidth}}}]{
                \texttt{\tiny{[GM$|$GM]}}
                {'player': 'Guesser', 'type': 'guess word is hallucination', 'utterance': 'mikroskop, laboratorium', 'message': "Guessed word 'laboratorium' was not listed, you can only guess words provided in the lists.", 'guess': 'laboratorium', 'previous_guesses': ['citron', 'oliv', 'peking', 'pilot', 'fågel fenix'], 'remaining_words': ['kraft', 'antarktis', 'hand', 'lager', 'pund', 'plan', 'bomb', 'ändra', 'obligation', 'modell', 'kung', 'mössa', 'våg', 'tokyo', 'triangel', 'papper', 'kort', 'stat', 'mikroskop', 'elev']}
            }
        }
    }
    & & \\ \\

\end{supertabular}
}

\end{document}
